\documentclass[11pt,a4paper]{article}

\usepackage[T1]{fontenc}
\usepackage[utf8]{inputenc}
\usepackage{lmodern}
\usepackage{amsmath,amssymb,amsthm}
\usepackage{booktabs,tabularx,array}
\usepackage{enumitem}
\usepackage[margin=2.5cm]{geometry}
\usepackage{microtype}
\usepackage[hidelinks,pdftitle={BRI-410: A 512-bit Tweakable Block Cipher with Key- and Tweak-Dependent Permutations},pdfauthor={Rinald Cobaj}]{hyperref}

\newtheorem{theorem}{Theorem}
\newtheorem{lemma}{Lemma}
\newtheorem{conjecture}{Conjecture}
\newtheorem*{corollary}{Corollary}

\newcommand{\ev}[1]{\textup{\textsf{[#1]}}}
\newcommand{\GF}{\mathrm{GF}}
\newcommand{\RK}{\mathit{RK}}
\newcommand{\shake}{\mathrm{SHAKE256}}
\newcommand{\str}[1]{\text{\texttt{"#1"}}}
\newcommand{\hx}[1]{\mathtt{#1}}
\newcommand{\Z}{\mathbb{Z}}

\setlist{itemsep=2pt,topsep=4pt}
\allowdisplaybreaks
\renewcommand{\arraystretch}{1.15}

\title{BRI-410: A 512-bit Tweakable Block Cipher\\with Key- and Tweak-Dependent Permutations}
\author{Rinald Cobaj\\Independent researcher\\\texttt{vlorarino@gmail.com}}
\date{Paper version 2.1, 5 November 2026}

\begin{document}
\maketitle

\begin{center}
\fbox{\begin{minipage}{0.92\textwidth}
\textbf{AI-use statement (IACR Policy on the Use of AI): Substantial AI contribution.}
The design ideas and decisions are the author's: the chess analogy and the 410-bit key, the per-block
``game played only once'', the change to 8-bit cells and 16 rounds, the 416-bit tag, the five
categories of analysis, and the choice of the schedule change in v1.1. An AI model (Claude, by
Anthropic), working under the author's direction, wrote the implementations, designed and ran the
experiments and computations (including the CP-SAT model), developed the proofs and the cryptanalysis
(including the weakness of Section~\ref{sec:weakness}), and wrote the text of this paper. The author
reviewed the results and is responsible for the content.
\end{minipage}}
\end{center}

\noindent\emph{This paper describes cipher version BRI-410 v1.1. BRI-410 is a study cipher: it has not been
analysed by independent cryptanalysts and must not be used to protect real data. Version 2.1 has the same
content as version 2 (27 October 2026), typeset with mathematical notation and with the AI-use statement above.}
Code, test vectors and all analysis scripts: \url{https://github.com/rinaldcobaj-cell/BRI-410}. Archive:
\url{https://doi.org/10.5281/zenodo.23693134}.

\begin{abstract}
BRI-410 is an experimental tweakable block cipher with a 512-bit block, a 410-bit key, a 128-bit tweak and
16 rounds. It comes with BRI-410 AE, an instantiation of the $\Theta$CB3 authenticated-encryption mode with a
416-bit tag. BRI-410 is a key-alternating substitution--permutation network on an $8\times 8$ array of bytes.
It uses an 8-bit S-box affine-equivalent to the AES S-box and an $8\times 8$ Cauchy MDS matrix. Its distinctive
feature is the permutation layer: in every round, the key and the tweak select a cell permutation from a family
of 32 ``knight-move'' permutations, all of which are diffusion-optimal in both directions. We prove that the
wide-trail bound of 81 active S-boxes over four rounds holds for every sequence of permutations from this
family. For individual differential characteristics and linear trails, this gives at least 324 active S-boxes
over 16 rounds. A CP-SAT model confirms the values for one to five rounds. This bound is not a proof of the
overall security of BRI-410. All round keys and permutations are derived from the key and the tweak.
Version 1.1 adds a Davies--Meyer feed-forward after tweak injection, intended to make the per-block derivation
one-way, because in version 1.0 three consecutive round keys of a single block revealed the master key in
0.1 seconds. That one-wayness is a conjecture, not a proven property. We document this weakness and the other
negative results of the design process, among them a predecessor with 4-bit cells that we broke over 5 rounds.
We report experiments in five categories of cryptanalysis, keep proofs, computations, experiments and
conjectures apart, and list open problems for independent analysis. We do not claim any practical advantage
over existing designs.
\end{abstract}

\section{Introduction}

The rules of chess are public, yet a serious game almost never repeats. BRI-410 starts from this analogy: the
algorithm is public, and the key, a ``game'', is secret. In addition, every encrypted block plays its own game,
derived from the key and a per-block tweak. The $8\times 8$ board fixes the state. The knight's move gives the
permutation layer. The key length comes from Allis's estimate of about $10^{123}$ games of chess~\cite{All94},
since $\log_2 10^{123} \approx 408.6$. \textbf{This choice is symbolic.} We do not claim that security above
256 bits has practical meaning, and our security target (Section~\ref{sec:claims}) is lower than $2^{410}$.

The aim of this paper is not to propose a replacement for existing ciphers. It describes BRI-410 completely
enough that independent cryptanalysts can analyse it, challenge it and, if possible, break it. It also records
every weakness we found during the design.

\paragraph{Contributions.}
\begin{enumerate}
\item A complete specification of BRI-410 v1.1, a 512-bit tweakable block cipher with a 410-bit key and
  128-bit tweak, and of BRI-410 AE (Sections~\ref{sec:spec} and~\ref{sec:ae}).
\item A proof that the two-round (9) and four-round (81) wide-trail bounds on individual differential
  characteristics and linear trails hold for every key- and tweak-dependent choice of the round permutations,
  both for the cipher and for the key schedule (Section~\ref{sec:struct}).
\item A record of the design process with its negative results. These include a useless parameter ($\alpha$),
  a key-dependent ``game'' that hides only one bit per round at row level, a 4-bit-cell predecessor broken over
  5 rounds, and an invertible key schedule in version 1.0, together with the change that addresses it
  (Sections~\ref{sec:weakness} and~\ref{sec:history}).
\item Experiments in five categories: differential, linear, algebraic, key schedule, and chosen
  plaintext/tweak (Section~\ref{sec:exp}). Open problems and concrete targets for analysis are in
  Section~\ref{sec:open}.
\end{enumerate}

\paragraph{Kinds of evidence.} Every statement is labelled:
\begin{itemize}
\item \ev{Proof}: a mathematical proof given in this paper;
\item \ev{Computation}: an exact computer computation (exhaustive check or solver);
\item \ev{Experiment}: a statistical or sample-based measurement;
\item \ev{Estimate}: a complexity estimate that was not implemented;
\item \ev{Conjecture}: a hypothesis we believe but cannot support with proof;
\item \ev{Open}: a question we did not study.
\end{itemize}

\section{Related work and positioning}

\paragraph{Tweakable block ciphers.} Liskov, Rivest and Wagner introduced tweakable block
ciphers~\cite{LRW02}. The TWEAKEY framework~\cite{JNP14} processes the key and the tweak with a light, linear
schedule. It underlies SKINNY~\cite{BJK+16} and Deoxys-BC~\cite{JNPS21}, which have 64- and 128-bit blocks and
extensive public analysis. Threefish~\cite{FLS+10}, the block cipher of Skein, is the closest design in size: it
has 256-, 512- and 1024-bit blocks and a 128-bit tweak, but it is an ARX design (addition, rotation, XOR) with
72 rounds for the 512-bit block.

\paragraph{Wide-block SPNs.} Rijndael supports 256-bit blocks~\cite{DR02}. Kalyna~\cite{OGK+15} is an SPN
standard with 512-bit blocks and keys, 8-bit S-boxes and $8\times 8$ MDS matrices over $\GF(2^8)$, but without
a tweak.

\paragraph{Key-dependent components and rekeying.} Blowfish~\cite{Sch93} and Twofish~\cite{SKW+98} derive
S-boxes from the key. Deriving a fresh key for every message or block is a known technique, analysed for example
in~\cite{AB00}. The Davies--Meyer construction from hash functions adds a feed-forward so that the output is
believed hard to invert when the underlying permutation is strong~\cite{PGV93}.

\begin{table}[ht]
\centering\small
\begin{tabularx}{\textwidth}{@{}lccccl>{\raggedright\arraybackslash}X@{}}
\toprule
Cipher & Block & Key & Tweak & Rounds & Structure & Tweak processing\\
\midrule
AES-256 & 128 & 256 & -- & 14 & SPN, 8-bit S-box, MDS & --\\
Kalyna-512/512 & 512 & 512 & -- & 18 & SPN, 8-bit S-boxes, MDS & --\\
Threefish-512 & 512 & 512 & 128 & 72 & ARX & added to subkeys\\
SKINNY, Deoxys-BC & 64 / 128 & \multicolumn{2}{c}{(tweakey)} & 14--56 & SPN & linear TWEAKEY schedule\\
\textbf{BRI-410 v1.1} & \textbf{512} & \textbf{410} & \textbf{128} & \textbf{16} & \textbf{SPN, 8-bit S-box, MDS} &
\textbf{full re-derivation of all round keys and permutations, intended to be one-way (Davies--Meyer, not proven)}\\
\bottomrule
\end{tabularx}
\caption{BRI-410 and related designs (sizes in bits).}
\end{table}

\paragraph{Positioning.} BRI-410 brings no gain in efficiency. Re-deriving the whole schedule for every block
is slower than TWEAKEY (Section~\ref{sec:perf}). What we study is a different design point: a permutation layer
that depends on the key and the tweak, chosen in every round from a family that is diffusion-optimal for every
choice. As a result, the classical wide-trail bounds hold uniformly over all keys and tweaks. To our knowledge
this exact combination has not been described before. Key-dependent components and tweak-dependent round keys
as such are well known.

\section{Specification of BRI-410 v1.1}\label{sec:spec}

\subsection{Notation and parameters}

\begin{itemize}
\item \emph{Block:} 512 bits $=$ 64 bytes. Byte $i$ is cell $(r,k) = (\lfloor i/8\rfloor,\, i \bmod 8)$ of an
  $8\times 8$ board; row $r$ consists of bytes $8r,\dots,8r+7$.
\item \emph{Key:} 410 bits, given as 52 bytes in little-endian order whose 6 top bits are $0$.
\item \emph{Tweak:} 128 bits (16 bytes).
\item \emph{Rounds:} $R = 16$.
\item \emph{Field:} $\GF(2^8) = \GF(2)[x]/(x^8 + x^4 + x^3 + x + 1)$. Bytes are identified with field elements;
  $\oplus$ denotes addition (XOR) and $\|$ concatenation.
\end{itemize}

\subsection{Components}

\paragraph{S-box $S$.} For $x \in \GF(2^8)$,
\[
  S(x) = A\cdot x^{-1} \oplus c, \qquad 0^{-1} := 0,
\]
where $A$ is an invertible $8\times 8$ matrix over $\GF(2)$ and $c$ is a byte. Row $i$ of $A$ is a byte $A_i$,
and bit $i$ of $A\cdot y$ is the parity of $A_i \wedge y$. The 9 bytes $A_0\|\cdots\|A_7\|c$ are the output of
$\shake(\str{BRI-410 S-box8 v1} \,\|\, n)$, where $n$ is the first 4-byte big-endian counter that gives an
invertible $A$ and an S-box with no $x$ such that $S(x) = x$ or $S(x) = x \oplus \hx{FF}$. The result is $n=9$.
The table of $S$ is in Appendix~\ref{app:const}.

\paragraph{Permutation family $\pi$.} Let
\[
  D = \bigl((1,2),\,(1,-2),\,(-1,2),\,(-1,-2)\bigr),
\]
in this order. For $d = (d_r, d_k) \in D$ and $\beta \in \{0,\dots,7\}$, the permutation $\pi_{d,\beta}$ of the
64 cells is
\[
  \pi_{d,\beta}(r,k) = \bigl(r + k\,d_r,\; k + k\,d_k + \beta\bigr) \bmod 8 .
\]
Cell $(r,k)$ makes $k$ knight jumps $(d_r,d_k)$ on a toroidal board, then the board shifts $\beta$ columns. This
gives 32 permutations, all distinct.

\paragraph{Linear layer $\theta$.} Each row, viewed as a vector in $\GF(2^8)^8$, is multiplied by the Cauchy
matrix
\[
  M = (M_{ij})_{0\le i,j\le 7}, \qquad M_{ij} = \frac{1}{x_i \oplus y_j},
\]
with $x = (0,1,\dots,7)$ and $y = (145, 234, 150, 206, 98, 226, 235, 94)$, the first 8 distinct bytes outside
$\{0,\dots,7\}$ of $\shake(\str{BRI-410 MDS8 v1})$. The matrices $M$ and $M^{-1}$ are in
Appendix~\ref{app:const}.

\subsection{Encryption and decryption}

Let $P$ be the plaintext, $\RK_0,\dots,\RK_{16}$ the round keys and $\pi_1,\dots,\pi_{16}$ the permutations
(Section~\ref{sec:schedule}). Encryption computes
\begin{align*}
  X &\gets P \oplus \RK_0, \\
  X &\gets \theta\bigl(\pi_i(S(X))\bigr) \oplus \RK_i \qquad \text{for } i = 1,\dots,16, \\
  C &\gets X .
\end{align*}
Here $S$ is applied to all 64 bytes, and $\pi_i$ moves the byte in cell $c$ of $S(X)$ to cell $\pi_i(c)$.
Decryption applies $\oplus\RK_i$, $\theta^{-1}$, $\pi_i^{-1}$, $S^{-1}$ for $i = 16,\dots,1$, and finally
$\oplus\RK_0$. All rounds are identical; the last round also applies $\theta$.

\subsection{Key and tweak schedule (``the game'')}\label{sec:schedule}

The schedule state $G$ has 80 bytes (10 rows of 8). For a byte constant $c$, one step is
\[
  F_c(G) = \sigma\bigl(\Theta(S(G))\bigr) \oplus c\, e_{79},
\]
where $S$ acts on all 80 bytes, $\Theta$ applies $\theta$ to each of the 10 rows, $\sigma$ moves byte $j$ to
position $9j \bmod 80$, and $c\,e_{79}$ adds $c$ to byte 79. Let $\mathit{pad}$ be the first 230 bits of
$\shake(\str{BRI-410/8 mbushja})$, read as a little-endian integer (Appendix~\ref{app:const}), and view the key
$K$ as a 410-bit integer. From $K$ and the tweak $T$, version 1.1 computes
\begin{align*}
  L &= K + 2^{410}\,\mathit{pad} &&\text{(80 bytes, little-endian)},\\
  K^* &= F_{\hx{F3}} \circ F_{\hx{F2}} \circ F_{\hx{F1}} \circ F_{\hx{F0}}(L)
      &&\text{(key only; the same for every tweak)},\\
  X &= K^* \oplus (T \,\|\, 0^{64}) &&\text{(the tweak is added to bytes } 0..15),\\
  G_0 &= F_{\hx{FB}} \circ \cdots \circ F_{\hx{F4}}(X) \oplus X &&\text{(8 steps and a Davies--Meyer feed-forward)},\\
  \RK_0 &= G_0[0..63] &&\text{(the whitening key)},
\end{align*}
and, for $i = 1,\dots,16$,
\[
  G_i = F_i(G_{i-1}), \qquad \RK_i = G_i[0..63], \qquad
  \pi_i = \pi_{D[b \bmod 4],\; \lfloor b/4\rfloor \bmod 8}, \quad b = G_i[64].
\]
So there are 17 round keys $\RK_0,\dots,\RK_{16}$ and 16 permutations $\pi_1,\dots,\pi_{16}$. Each step
evaluates 80 S-boxes; the cost of the schedule is given in Table~\ref{tab:cost}.

\begin{table}[ht]
\centering\small
\begin{tabular}{@{}lrr@{}}
\toprule
Process & Steps & S-boxes\\
\midrule
Key preparation $K^*$, once per key and shared by all tweaks & 4 & 320\\
Tweak mixing, per tweak & 8 & 640\\
Round-key generation, per tweak & 16 & 1\,280\\
Schedule for one tweak, without the shared preparation & 24 & 1\,920\\
Schedule for the first tweak under a key, including the preparation & 28 & 2\,240\\
\bottomrule
\end{tabular}
\caption{Cost of the key and tweak schedule of v1.1.}\label{tab:cost}
\end{table}

With the 1\,024 S-boxes of the 16 rounds, the first block under a key costs 3\,264 S-boxes, and every further
block, which reuses $K^*$, costs 2\,944. The feed-forward is a design measure meant to make the map from $G_0$
back to $X$ hard to invert. \textbf{We do not prove that it is one-way} (Conjecture~\ref{conj:isolation} in
Section~\ref{sec:claims}).

\paragraph{Version 1.0} (Sections~\ref{sec:exp} to~\ref{sec:attacks}, history) used 4 steps after the tweak and
no feed-forward, $G_0 = F_{\hx{F7}} \circ \cdots \circ F_{\hx{F4}}(X)$, which is $4 + 16 = 20$ steps per tweak
after the shared preparation, or 24 including it. The round function is the same in both versions.

\section{BRI-410 AE}\label{sec:ae}

BRI-410 AE instantiates $\Theta$CB3~\cite{KR11} with BRI-410/8 v1.1 as the tweakable block cipher $E$. Each call
uses the tweak
\[
  N \,\|\, \mathrm{LE}_{32}\bigl(16\cdot \mathit{index} + \mathit{type}\bigr),
\]
that is, a 12-byte nonce $N$ followed by 4 bytes in little-endian order, with $\mathit{index} < 2^{28}$. We
write $E^{N,i,t}$ for $E$ under this tweak. The types are: 0 message block; 1 mask of the final partial block;
2 tag without a partial block; 3 tag with a partial block; 4 associated-data block; 5 final partial
associated-data block. Associated data uses $N = 0^{96}$. Let $M = M_1\cdots M_m M_*$ and $A = A_1\cdots A_a A_*$
with full 64-byte blocks $M_i, A_i$ and $|M_*|, |A_*| < 64$ bytes, and let $\mathit{pad}(x) = x \,\|\, \hx{80}
\,\|\, \hx{00}\cdots\hx{00}$. With $\Sigma = M_1 \oplus \cdots \oplus M_m$:
\begin{align*}
  C_i &= E^{N,i,0}(M_i), \qquad i = 1,\dots,m,\\
  \text{if } M_* \ne \varepsilon:\quad C_* &= M_* \oplus \mathrm{trunc}_{|M_*|}\bigl(E^{N,m,1}(0^{512})\bigr),
     \qquad F = E^{N,m,3}\bigl(\Sigma \oplus \mathit{pad}(M_*)\bigr),\\
  \text{otherwise:}\quad F &= E^{N,m,2}(\Sigma),\\
  \mathit{Auth} &= \bigoplus_{i=1}^{a} E^{0,i,4}(A_i)\; \bigl[\oplus\, E^{0,a,5}\bigl(\mathit{pad}(A_*)\bigr)
     \text{ if } A_* \ne \varepsilon\bigr],\\
  \mathit{Tag} &= \mathrm{trunc}_{416}(F \oplus \mathit{Auth}),\\
  \mathit{Packet} &= N \,\|\, C_1 \cdots C_m C_* \,\|\, \mathit{Tag}.
\end{align*}
Decryption returns nothing unless the tag verifies, and the comparison is constant-time. The correspondence with
$\Theta$CB3 was checked by the author against the description of the mode, not by an independent party.
\ev{Open}

\section{Structural results}\label{sec:struct}

\begin{lemma}[permutation family]\label{lem:perm}
\ev{Proof} Each $\pi_{d,\beta}$ is a bijection, and it is diffusion-optimal in both directions: the 8 cells of
any row go to 8 distinct rows, and the 8 cells that land in any row come from 8 distinct rows.
\end{lemma}

\begin{proof}
The column map $k \mapsto k(1 + d_k) + \beta$ is a bijection of $\Z_8$, since $1 + d_k \in \{3, -1\}$ is odd.
Given the column, the row map $r \mapsto r + k\,d_r$ is a bijection, so $\pi_{d,\beta}$ is a bijection. For a
fixed source row $r$, the destination rows $r + k\,d_r$ $(k = 0,\dots,7)$ are distinct because $d_r = \pm 1$ is
invertible modulo 8. For a destination row $q$, every column $k$ occurs exactly once among its cells, because the
column map is a bijection. The source row of the cell with column $k$ is $q - k\,d_r$, and these rows are
distinct.
\end{proof}

\begin{lemma}[the row shift $\alpha$ adds no structure]\label{lem:alpha}
\ev{Proof} Let $T_\alpha$ be the row rotation $(r,k) \mapsto (r+\alpha, k)$. Then $T_\alpha$ commutes with $S$,
with every $\pi_{d,\beta}$ and with $\theta$. Hence a round-dependent $T_\alpha$ inside the round only adds a
final, key-dependent rotation of rows.
\end{lemma}

\begin{proof}
$S$ acts cell by cell. For $\pi_{d,\beta}$, both $\pi_{d,\beta}\circ T_\alpha$ and $T_\alpha \circ \pi_{d,\beta}$
map $(r,k)$ to $\bigl(r + \alpha + k\,d_r,\; k(1+d_k) + \beta\bigr)$. The layer $\theta$ applies the same matrix
to every row. Therefore $T_\alpha$ can be moved to the end of the cipher, replacing $\RK_i$ with
$T(\RK_i)$. We therefore removed $\alpha$; in an earlier version, measurements found 1\,792 pairs of ``game''
variants that were equivalent up to $\alpha$.
\end{proof}

\begin{lemma}[MDS]\label{lem:mds}
\ev{Proof + Computation} $M$ is MDS, and its differential and linear branch numbers are both 9.
\end{lemma}

\begin{proof}
Every square submatrix of a Cauchy matrix is a Cauchy matrix, and its determinant is non-zero when the $x_i$ are
distinct, the $y_j$ are distinct and $x_i \ne y_j$~\cite{MS77}. Hence $M$ is MDS and its differential branch
number is 9. Since the dual of an MDS code is MDS, the linear branch number is also 9. We also checked all
12\,869 square submatrices of $M$ and of $M^{-1}$ by computer: none is singular. $M$ has no column symmetry other
than the identity.
\end{proof}

\begin{theorem}[bounds for every permutation sequence]\label{thm:bounds}
\ev{Proof} Consider any differential (or linear) trail through $r$ rounds of BRI-410 with a non-zero input, for
any key and any tweak. Then, for every sequence of permutations $\pi_i$ from the family,
\begin{enumerate}[label=(\alph*)]
\item any two consecutive rounds contain at least 9 active S-boxes, and
\item any four consecutive rounds contain at least 81 active S-boxes.
\end{enumerate}
\end{theorem}

\begin{proof}
Let $a_t$ be the number of active bytes at the input of the S-layer of round $t$, and let $\rho(\cdot)$ denote
the number of active rows of a state.

(a) $\theta_t$ acts row by row with branch number 9: for each active row at its input, the active input bytes
plus the active output bytes number at least 9. The input bytes of $\theta_t$ are the S-outputs of round $t$,
moved by $\pi_t$. Its output bytes are the S-inputs of round $t+1$, since adding the key does not change
activity. Summing over the rows gives
\[
  a_t + a_{t+1} \;\ge\; 9\,\rho(\text{input of } \theta_t) \;\ge\; 9 .
\]

(b) Take rounds 1 to 4 and write $\rho_t = \rho(\text{input of } \theta_t)$. By (a),
$a_1 + a_2 \ge 9\rho_1$ and $a_3 + a_4 \ge 9\rho_3$. Let $q$ be any active row at the input of $\theta_2$. It has
$w_{\mathrm{in}}$ active bytes, and its image row at the output of $\theta_2$ has $w_{\mathrm{out}}$ active bytes,
with $w_{\mathrm{in}} + w_{\mathrm{out}} \ge 9$. The $w_{\mathrm{in}}$ bytes come, through $\pi_2$, from
$w_{\mathrm{in}}$ distinct rows of the state at the input of round 2 (Lemma~\ref{lem:perm}, inverse direction).
Each of these rows is an output row of $\theta_1$, so it is also active at the input of $\theta_1$; hence
$\rho_1 \ge w_{\mathrm{in}}$. The $w_{\mathrm{out}}$ bytes of row $q$ are S-inputs of round 3, and $\pi_3$ sends
them to $w_{\mathrm{out}}$ distinct rows (Lemma~\ref{lem:perm}, forward direction); hence
$\rho_3 \ge w_{\mathrm{out}}$. Therefore
\[
  a_1 + a_2 + a_3 + a_4 \;\ge\; 9(\rho_1 + \rho_3) \;\ge\; 9(w_{\mathrm{in}} + w_{\mathrm{out}}) \;\ge\; 81 .
\]
The argument for linear trails is the same, with masks propagating through the transpose of $\theta$, which also
has branch number 9 (Lemma~\ref{lem:mds}).
\end{proof}

\begin{corollary}
$4k$ rounds contain at least $81k$ active S-boxes; in particular, 16 rounds contain at least 324.
\end{corollary}

\begin{theorem}[key and tweak schedule]\label{thm:schedule}
\ev{Proof + Computation} The permutation $j \mapsto 9j \bmod 80$ of the schedule sends the 8 bytes of each of the
10 rows to 8 distinct rows, and receives each row from 8 distinct rows \ev{Computation}. The proof of
Theorem~\ref{thm:bounds} therefore applies to the steps $F_c$: any non-zero difference in the schedule state, in
particular a difference in the tweak, activates at least 81 S-boxes in 4 steps and at least 162 in the 8 steps of
v1.1 that precede $\RK_0$.
\end{theorem}

\paragraph{What these bounds guarantee, and what they do not.}
\begin{itemize}
\item They \textbf{guarantee} that every single differential characteristic, or linear trail, activates at least
  the stated number of S-boxes, for every key, every tweak and every choice of permutations. With the maximal
  S-box probability $2^{-6}$, any single 4-round trail has probability at most $(2^{-6})^{81} = 2^{-486}$ and any
  16-round trail at most $2^{-1944}$, \textbf{under the usual assumption of independent round keys}.
\item They \textbf{do not guarantee} anything about (i) differentials, where many trails add up, (ii) linear
  hulls, (iii) the actual dependencies between round keys produced by the schedule, or (iv) attacks of other
  kinds: integral, algebraic, boomerang, impossible-differential, or related-key beyond the schedule bound.
\item In particular, Theorems~\ref{thm:bounds} and~\ref{thm:schedule} are \textbf{not a proof of the security of
  BRI-410 as a whole}.
\end{itemize}

\paragraph{Computation. \ev{Computation}} A CP-SAT model at the level of active bytes, in which the adversary may
choose the row structure of $\pi_i$ in every round (Appendix~\ref{app:cpsat}), returned the minimum number of
active S-boxes: 1, 9, 17, 81 and 82 for 1 to 5 rounds (proven optimal), and 1, 9, 17, 81, 107 and 114 for 1 to 6
schedule steps. For 6, 7 and 8 rounds it found trails with exactly 90, 98 and 162 active S-boxes, which equal
the bounds that follow from Theorem~\ref{thm:bounds}. These values were computed once (OR-Tools 9.15,
22 October 2026). Later runs on the author's machine were blocked by an operating-system application-control
policy, so an independent re-run is welcome. The proofs above do not depend on this computation.

\section{Experiments in five categories}\label{sec:exp}

All experiments use fixed random seeds and are in the repository. Experiments marked ``v1.0'' were run before
the schedule change. Experiments that only involve the round function under a fixed tweak (D, L, A) do not
depend on how the round keys are derived. Experiments K and T were repeated for v1.1
(Section~\ref{sec:weakness}).

\subsection{Differential (D)}

\begin{center}\small
\begin{tabularx}{\textwidth}{@{}>{\raggedright\arraybackslash}X>{\raggedright\arraybackslash}Xl@{}}
\toprule
Measurement & Result & \\
\midrule
S-box difference table ($255 \times 256$ entries) & $32\,895\times 0$, $32\,130\times 2$, $255\times 4$;
  maximal probability $2^{-6}$ & \ev{Computation}\\
One round: most frequent output difference ($3\cdot 2^{14}$ pairs) & frequency 0.0161 / 0.0157 / 0.0142, against
  $4/256 = 0.0156$ predicted; 127 distinct outputs, as predicted & \ev{Experiment}\\
One active input byte: number of output bytes with zero difference (300 pairs; random: 0.25) & 1 round: 56;
  2 rounds: 0 in every pair; 3 rounds: 0.29 & \ev{Experiment}\\
\bottomrule
\end{tabularx}
\end{center}

After 2 rounds every byte is always active, which distinguishes 2 rounds from a random permutation, as for AES.
We saw no deviation after 3 rounds.

\subsection{Linear (L)}

\begin{center}\small
\begin{tabularx}{\textwidth}{@{}>{\raggedright\arraybackslash}X>{\raggedright\arraybackslash}Xl@{}}
\toprule
Measurement & Result & \\
\midrule
S-box: maximal correlation & $32/256 = 2^{-3}$ & \ev{Computation}\\
One round, exact: only byte 0 varies (256 inputs), output byte 0 & maximal correlation 0.125, equal to the
  S-box; a byte of another row does not change at all & \ev{Computation}\\
\bottomrule
\end{tabularx}
\end{center}

\paragraph{Negative methodological result.} We first ran a statistical test that correlates one input byte with
one output byte over $2^{18}$ random plaintexts. It is \textbf{uninformative by construction}: with all other
input bytes random, an output byte is the XOR of 8 S-box outputs from 8 different input bytes, so the
correlation is zero even after one round. We do not use it as evidence. Linear trails over two or more rounds
have correlation at most $2^{-27}$ and cannot be measured experimentally. \textbf{The linear security argument
rests on Theorem~\ref{thm:bounds}.}

\subsection{Algebraic (A)}

\begin{center}\small
\begin{tabularx}{\textwidth}{@{}>{\raggedright\arraybackslash}X>{\raggedright\arraybackslash}Xl@{}}
\toprule
Measurement & Result & \\
\midrule
Linearly independent quadratic equations relating input and output of $S$ & 39, equal to the AES S-box with the
  same code & \ev{Computation}\\
Linear structures of $S$ & 0 & \ev{Computation}\\
Size of a polynomial system that expresses one block in terms of the master key, v1.1 & 3\,264 S-boxes (1\,024 in
  the cipher and 2\,240 in the schedule, the 4 key-only steps included; Section~\ref{sec:schedule}):
  $\approx 127\,000$ quadratic equations (39 per S-box) in $\approx 26\,500$ variables (8 per S-box output plus
  410 key bits). Further blocks under the same key add 2\,944 S-boxes each. This is a count, not an attack &
  \ev{Computation}\\
\bottomrule
\end{tabularx}
\end{center}

\paragraph{Cube tests \ev{Experiment}.} We computed the XOR-sum of the outputs over all values of $d$ chosen
plaintext bits under a fixed tweak, and counted the output bits whose sum is zero (random: about 256 of 512).

\begin{center}\small
\begin{tabular}{@{}lccc@{}}
\toprule
Cube & 2 rounds & 3 rounds & 4 rounds\\
\midrule
8 bits in one byte & 512, 512 & 512, 512 & 241, 264\\
8 bits in 8 different bytes & 512, 512 & 236, 263 & 247, 255\\
12 bits in 12 bytes & 512, 512 & 259, 234 & 258, 248\\
16 bits in 16 bytes & 512, 512 & 512, 257 & 258, 232\\
\bottomrule
\end{tabular}
\end{center}

\paragraph{Explanation of the 3-round zero sums \ev{Proof}.} In round 1 each S-box sees at most one variable bit,
so its output is affine in that bit. If $\pi_1$ puts at most $k$ cube bits in each row, the degree after round 2
is at most $k$ and after round 3 at most $7k$. With $8k$ cube bits in total and $7k < 8k$, the sum is zero. We
observed no zero sums over 4 rounds.

\subsection{Key schedule (K)}

\begin{center}\small
\begin{tabularx}{\textwidth}{@{}>{\raggedright\arraybackslash}X>{\raggedright\arraybackslash}Xl@{}}
\toprule
Measurement & Result & \\
\midrule
One key bit flipped: bits changed in $\RK_0, \RK_1, \RK_8, \RK_{16}$ (v1.1, 820 trials) & 50.1\%, 50.1\%,
  50.1\%, 50.0\% & \ev{Experiment}\\
Permutation changes when one key bit flips (v1.0) & 97.0\%, against $1 - 1/32 = 96.9\%$ expected &
  \ev{Experiment}\\
Structured keys (zero, all ones, single bit, $\hx{55}\ldots$, counting; v1.0) & round keys 50--51\% ones; 3\,000
  output blocks each: monobit $z \le 1.98$, byte $\chi^2$ 218--275 (99\% threshold 310.5) & \ev{Experiment}\\
\textbf{Recovery of the master key from round keys} & v1.0: broken. v1.1: blocked for the direct method
  (Section~\ref{sec:weakness}) & \ev{Experiment}\\
\bottomrule
\end{tabularx}
\end{center}

\subsection{Chosen plaintext and chosen tweak (T)}

\begin{center}\small
\begin{tabularx}{\textwidth}{@{}>{\raggedright\arraybackslash}X>{\raggedright\arraybackslash}Xl@{}}
\toprule
Measurement & Result (random: 0.25 bytes) & \\
\midrule
Tweak integral, v1.1 (one tweak byte takes all 256 values, plaintext fixed) after 1 / 2 / 16 rounds & 0--1 / 0--1
  / 0--1 of 64 bytes balanced & \ev{Experiment}\\
Tweak integral, v1.1 (two tweak bytes, 65\,536 tweaks), 1 round & 0 of 64 & \ev{Experiment}\\
Joint difference in plaintext and tweak chosen to cancel the difference of $\RK_0$, v1.1, rounds 1 / 2 / 3 &
  0.23 / 0.28 / 0.26 bytes with zero difference & \ev{Experiment}\\
Same with \textbf{no} mixing after the tweak (weakened variant), rounds 1 / 2 / 3 & 56.5 / 0.59 / 0.27 &
  \ev{Experiment}\\
16 rounds, tweak $= 0, 1, 2, \dots$, plaintext 0, 10\,000 blocks (v1.1) & monobit $z = -0.74$, runs $z = 0.78$,
  byte $\chi^2 = 261$, nibble $\chi^2 = 6.1$, all below the 99\% thresholds & \ev{Experiment}\\
\bottomrule
\end{tabularx}
\end{center}

Statistical tests detect only gross defects and are not evidence of security.

\section{The key-schedule weakness of v1.0 and the change in v1.1}\label{sec:weakness}

\subsection{The weakness}\label{sec:weakness71}

In v1.0 every schedule step is a bijection, and the steps after the tweak have no feed-forward. \textbf{Given
three consecutive round keys $\RK_i, \RK_{i+1}, \RK_{i+2}$ of one block, and its tweak, the master key follows
immediately.} \ev{Experiment}
\begin{enumerate}
\item $\RK_i$ gives bytes $0..63$ of $G_i$. These determine the outputs of $\theta$ on rows $0..7$ of step $i+1$,
  so 13 of the 16 bytes $64..79$ of $G_{i+1}$ are known. Bytes 70, 71 and 79 remain unknown.
\item $\RK_{i+2}$ gives 6 byte equations on row 8 and 7 on row 9 of $G_{i+1}$. An exhaustive search over $2^{16}$
  and $2^8$ values gives a unique $G_{i+1}$.
\item The steps $F_i$ are inverted back to $G_0$. The 4 tweak-mixing steps are undone, the tweak is removed, and
  the 4 key-only steps are undone. The 230 known padding bits confirm the key.
\end{enumerate}
In 3 out of 3 trials the master key was recovered in about 0.1\,s, with a single candidate. With \textbf{two}
consecutive round keys, 3 bytes remain unknown: $2^{24}$ candidates, each verifiable through the padding bits
\ev{Estimate}. Once the key is known, the round keys of \textbf{every} other tweak follow. The per-block ``game''
therefore did not isolate blocks. This weakness does not break 16 rounds, since no method is known to obtain round
keys there. It does mean that, in v1.0, any attack recovering two or three consecutive round keys of one block
becomes a full break.

\subsection{The change}

Version 1.1 adds a Davies--Meyer feed-forward after the tweak injection and doubles the number of steps there from
4 to 8: $G_0 = F(X) \oplus X$, with $F$ made of 8 steps and $X = K^* \oplus (T \,\|\, 0^{64})$. The state $G_0$ of
an attacked block can still be reconstructed. But going from $G_0$ back to $X$ now requires inverting the map
$X \mapsto F(X) \oplus X$, a preimage problem on 640 bits. The step-by-step inversion of the previous subsection
no longer applies, because the feed-forward mixes the input into the output. \textbf{Whether the map can be
inverted by other means is not proven either way}: the feed-forward is a design measure, and its one-wayness is
Conjecture~\ref{conj:isolation}. The round function is unchanged.

\subsection{Evaluation}

\begin{center}\small
\begin{tabular}{@{}lccc@{}}
\toprule
Attack of Section~\ref{sec:weakness71} & $G_0$ of the & Master key & Round keys of\\
(3 trials each) & attacked block & recovered & another tweak\\
 & reconstructed & & predicted\\
\midrule
v1.0 & 3/3 & \textbf{3/3} & \textbf{3/3}\\
v1.1 & 3/3 & \textbf{0/3} & \textbf{0/3}\\
\bottomrule
\end{tabular}
\end{center}

Cost: v1.1 is 13\% slower than v1.0 in our implementation (2\,295 against 2\,651 blocks per second, with a new
tweak for every block and the key preparation $K^*$ shared).

\textbf{This is not a proof.} The three failures show only that the direct inversion of Section~\ref{sec:weakness71} no longer
works. They do not show that $F(X) \oplus X$ cannot be inverted by other methods, such as meet-in-the-middle or
rebound techniques of the kind applied to AES-based hashing modes~\cite{Sas11}. Our hypothesis is stated as
Conjecture~\ref{conj:isolation}.

\section{Reduced-round attacks (fixed tweak)}\label{sec:attacks}

\begin{lemma}[3-round integral distinguisher]\label{lem:integral}
\ev{Proof + Experiment} Under a fixed tweak, take 256 plaintexts in which one byte takes all values and the rest is
constant. After 3 rounds, every one of the 64 output bytes sums to zero.
\end{lemma}

\begin{proof}
After round 1, the row that receives the active byte contains 8 bytes of the form $c \oplus m\cdot S(v)$ with
$m \ne 0$, so each takes all 256 values. In round 2, $\pi_2$ sends these 8 bytes to 8 distinct rows
(Lemma~\ref{lem:perm}). Each row then has one byte that takes all values and 7 constant bytes, so after $\theta$ all
64 bytes take all values. In round 3, $\theta$ outputs XOR combinations of such bytes, so every output byte sums to
zero. Adding the key keeps this property.
\end{proof}

Measured: 10 out of 10 trials were balanced after 2 and 3 rounds. After 4 rounds, 0--1 of 64 bytes were balanced,
and 0--1 also with two aligned bytes ($2^{16}$ texts). With a fresh tweak per plaintext the property disappears
\ev{Experiment}.

\paragraph{4-round key recovery \ev{Experiment}.} Lemma~\ref{lem:integral}, together with guessing each byte of
$\theta^{-1}(\RK_4)$ separately, recovered $\RK_4$ (512 bits) in 5 out of 5 trials with 512--768 chosen
plaintexts and in under one second. This works because
\[
  \theta^{-1}(C) = \pi_4\bigl(S(X_3)\bigr) \oplus \theta^{-1}(\RK_4)
\]
and every byte of $X_3$ is balanced. \textbf{This attack recovers only $\RK_4$, not the master key.} Under v1.0,
the master key would follow from three consecutive round keys of one block, or from two with about $2^{24}$ work
(Section~\ref{sec:weakness}). Removing round 4 with the recovered $\RK_4$ and attacking the remaining rounds in
the same way might also give $\RK_3$ and $\RK_2$, but we did not implement this. Whether a 4-round master-key
recovery on v1.0 is practical therefore remains \ev{Open}. Under v1.1, round keys of one tweak do not lead to the
master key by the direct method of Section~\ref{sec:weakness}.

\paragraph{5 rounds \ev{Estimate}.} Adding one more round requires guessing 8 bytes of $\theta^{-1}(\RK_5)$ and
one byte of $\theta^{-1}(\RK_4)$ per target byte: about $2^{72}$ guesses with $2^8$ texts, plus the unknown row
structure of $\pi_4$ and $\pi_5$. Done naively, this is about $2^{80}$ S-box evaluations. Partial-sum techniques
may reduce it. We did not implement this attack, and the estimate may be wrong in either direction.

\paragraph{Margin.} The best attack we executed covers 4 of 16 rounds. Our rule of thumb
$R \ge 2\max(R_1, R_2, R_3)$, a heuristic and not a proof, is satisfied. Here $R_1 = 2$ rounds to full diffusion,
$R_2 = 6$ rounds until the best trail is below $2^{-512}$ (by Theorem~\ref{thm:bounds}), and $R_3 = 4$--$5$ rounds
attacked, so $2\cdot 6 = 12 \le 16$.

\section{Security claims and conjectures}\label{sec:claims}

\paragraph{What is proven:} Lemmas~\ref{lem:perm}--\ref{lem:integral}, Theorems~\ref{thm:bounds}
and~\ref{thm:schedule} and their corollaries, as stated, including their limits (Section~\ref{sec:struct}).

\begin{conjecture}[security target]\label{conj:target}
In the single-key setting, with chosen plaintexts, ciphertexts and tweaks (repeated tweaks included), 16-round
BRI-410 v1.1 has no distinguishing or key-recovery attack with time below $2^{256}$ and data below $2^{128}$ blocks
per key.
\end{conjecture}

This is a research hypothesis and an invitation to analysis, \textbf{not a demonstrated result}. We deliberately set
this target far below the key length, and the key length remains 410 bits. The 410-bit key is a symbolic choice, and
our analysis covers only the attack classes of Sections~\ref{sec:struct}--\ref{sec:attacks}.

\begin{conjecture}[tweak isolation]\label{conj:isolation}
The map $X \mapsto F(X) \oplus X$ with its 8 steps is one-way in practice. As a result, knowing all round keys and
permutations for one tweak gives no practical advantage in computing those of another tweak or the master key.
\end{conjecture}

This is a hypothesis supported only by the design rationale and by the failure of the direct attack
(Section~\ref{sec:weakness}).

\paragraph{BRI-410 AE.} If Conjecture~\ref{conj:target} holds, the security theorem of $\Theta$CB3~\cite{KR11}
gives confidentiality and authenticity against nonce-respecting adversaries. The forgery probability per attempt is
about $2^{-416}$ plus the advantage against the cipher. Limits: messages of at most $2^{28} - 1$ blocks. With random
96-bit nonces, at most $2^{32}$ messages per key, which keeps the probability of a repeated nonce around $2^{-33}$.

\paragraph{Not claimed:}
\begin{itemize}
\item security against related keys;
\item security of BRI-410 inside a hash function (known or chosen key);
\item resistance to nonce reuse in BRI-410 AE (a repeated nonce reveals equal blocks);
\item protection of the reference code against timing and other side channels (it is not constant-time).
\end{itemize}

\section{Design history and negative results}\label{sec:history}

\begin{center}\small
\begin{tabularx}{\textwidth}{@{}>{\raggedright\arraybackslash}p{2.6cm}>{\raggedright\arraybackslash}X%
>{\raggedright\arraybackslash}X>{\raggedright\arraybackslash}p{3.2cm}@{}}
\toprule
Version & Description & What we found & Action\\
\midrule
BRI-410/4, first design & 4-bit cells, 256-bit block, ``game'' ($\pi$) derived from the key only, $\alpha$
  present & $\alpha$ commutes with the round (Lemma~\ref{lem:alpha}): 1\,792 equivalent pairs & $\alpha$ removed\\
BRI-410/4, game per key & as above & the game hides only 1 bit per round at row level (the destination row depends
  only on $d_r$) & documented\\
BRI-410/4, game per key & as above & integral distinguisher over \textbf{4 rounds with $2^{12}$ texts},
  \textbf{5-round key recovery in under 1\,s} & game per block (tweak) introduced: the same attack failed 0/3\\
BRI-410/4 with tweak & 32 rounds & the 4-bit S-box (degree 3) lets the integral property survive one round more
  than in AES-like designs & replaced by 8-bit cells\\
BRI-410/8 v1.0 & 8-bit cells, 512-bit block, 16 rounds & invertible schedule: 3 round keys $\to$ master key in
  0.1\,s (Section~\ref{sec:weakness}) & v1.1 (feed-forward, 8 steps)\\
BRI-410/8 v1.1 & this paper & -- & open for analysis\\
\bottomrule
\end{tabularx}
\end{center}

\paragraph{Errors corrected from version 1 of this paper.}
\begin{itemize}
\item It stated a security claim at the level of $2^{410}$ that our analysis did not support.
\item It called the second implementation ``independent'', although it was written separately by the same author.
\item It called the CP-SAT values ``tight'' without saying that they are tight only within the model.
\item It reported ``49.9\%'' for a value whose measurement is 49.8\%.
\item It described cipher version v1.0.
\end{itemize}

\section{Reproducibility and performance}\label{sec:perf}

The repository contains:
\begin{itemize}
\item the research implementation (Python, standard library);
\item a separate single-file reference implementation in English, which agrees bit for bit with the research
  implementation;
\item test vectors for v1.0 and v1.1;
\item 47 unit tests;
\item every experiment script, with fixed seeds.
\end{itemize}
The CP-SAT model requires OR-Tools. In pure Python the implementation runs at about 2\,300 blocks per second with a
new tweak for every block and the key preparation $K^*$ shared, and BRI-410 AE at about 150\,KB/s. A constant-time
implementation in C is an open task.

\section{Open problems and targets for analysis}\label{sec:open}

\begin{enumerate}
\item \textbf{Differential and linear:} clustering of differentials, linear hulls, boomerang and
  impossible-differential attacks; re-running the CP-SAT bounds independently.
\item \textbf{Algebraic and integral:} division-property analysis for 4 or more rounds, and algebraic attacks on
  reduced rounds.
\item \textbf{Key schedule:} preimages of $X \mapsto F(X) \oplus X$ with 8 steps (Conjecture~\ref{conj:isolation});
  related-key attacks; whether reduced-round attacks lead to the master key.
\item \textbf{Tweak:} related-tweak attacks that combine tweak and plaintext differences, including pairs whose
  permutations differ.
\item \textbf{Mode and implementation:} an independent check of the $\Theta$CB3 instantiation, and constant-time
  implementations.
\end{enumerate}

\paragraph{Targets.} Distinguishers for 5, 6, 7 or 8 rounds; recovery of round keys for 5 to 8 rounds; recovery of
the master key, or of the round keys of a second tweak, from those of one tweak, in either the fixed-tweak or the
chosen-tweak setting. The author will acknowledge every result publicly.

\section*{Intellectual property}

BRI-410 is not patented, and the author does not intend to apply for any patent on it. The specification, the
implementations, the test vectors and the analysis scripts are dedicated to the public domain under CC0~1.0.

\section*{Use of AI}

See the AI-use statement on the first page. Every result can be reproduced with the published code.

\begin{thebibliography}{MWGP11}
\bibitem[AB00]{AB00} M.~Abdalla, M.~Bellare. Increasing the lifetime of a key: a comparative analysis of the
  security of re-keying techniques. \emph{ASIACRYPT 2000}.
\bibitem[All94]{All94} L.~V.~Allis. \emph{Searching for Solutions in Games and Artificial Intelligence.} PhD
  thesis, University of Limburg, Maastricht, 1994.
\bibitem[BJK+16]{BJK+16} C.~Beierle, J.~Jean, S.~K\"olbl, G.~Leander, A.~Moradi, T.~Peyrin, Y.~Sasaki,
  P.~Sasdrich, S.~M.~Sim. The SKINNY family of block ciphers and its low-latency variant MANTIS.
  \emph{CRYPTO 2016}.
\bibitem[BS91]{BS91} E.~Biham, A.~Shamir. Differential cryptanalysis of DES-like cryptosystems. \emph{Journal of
  Cryptology}, 1991.
\bibitem[DKR97]{DKR97} J.~Daemen, L.~Knudsen, V.~Rijmen. The block cipher Square. \emph{FSE 1997}.
\bibitem[DR02]{DR02} J.~Daemen, V.~Rijmen. \emph{The Design of Rijndael.} Springer, 2002.
\bibitem[FLS+10]{FLS+10} N.~Ferguson, S.~Lucks, B.~Schneier, D.~Whiting, M.~Bellare, T.~Kohno, J.~Callas,
  J.~Walker. The Skein hash function family. SHA-3 submission, version 1.3, 2010.
\bibitem[JNP14]{JNP14} J.~Jean, I.~Nikoli\'c, T.~Peyrin. Tweaks and keys for block ciphers: the TWEAKEY framework.
  \emph{ASIACRYPT 2014}.
\bibitem[JNPS21]{JNPS21} J.~Jean, I.~Nikoli\'c, T.~Peyrin, Y.~Seurin. The Deoxys AEAD family. \emph{Journal of
  Cryptology}, 2021.
\bibitem[KR11]{KR11} T.~Krovetz, P.~Rogaway. The software performance of authenticated-encryption modes.
  \emph{FSE 2011}.
\bibitem[LRW02]{LRW02} M.~Liskov, R.~Rivest, D.~Wagner. Tweakable block ciphers. \emph{CRYPTO 2002}.
\bibitem[Mat93]{Mat93} M.~Matsui. Linear cryptanalysis method for DES cipher. \emph{EUROCRYPT 1993}.
\bibitem[MS77]{MS77} F.~J.~MacWilliams, N.~J.~A.~Sloane. \emph{The Theory of Error-Correcting Codes.}
  North-Holland, 1977.
\bibitem[MWGP11]{MWGP11} N.~Mouha, Q.~Wang, D.~Gu, B.~Preneel. Differential and linear cryptanalysis using
  mixed-integer linear programming. \emph{Inscrypt 2011}.
\bibitem[OGK+15]{OGK+15} R.~Oliynykov, I.~Gorbenko, O.~Kazymyrov et al. A new encryption standard of Ukraine: the
  Kalyna block cipher. \emph{Cryptology ePrint Archive}, 2015/650.
\bibitem[PGV93]{PGV93} B.~Preneel, R.~Govaerts, J.~Vandewalle. Hash functions based on block ciphers: a synthetic
  approach. \emph{CRYPTO 1993}.
\bibitem[Sas11]{Sas11} Y.~Sasaki. Meet-in-the-middle preimage attacks on AES hashing modes and an application to
  Whirlpool. \emph{FSE 2011}.
\bibitem[Sch93]{Sch93} B.~Schneier. Description of a new variable-length key, 64-bit block cipher (Blowfish).
  \emph{FSE 1993}.
\bibitem[SKW+98]{SKW+98} B.~Schneier, J.~Kelsey, D.~Whiting, D.~Wagner, C.~Hall, N.~Ferguson. Twofish: a 128-bit
  block cipher. AES submission, 1998.
\bibitem[Sha50]{Sha50} C.~E.~Shannon. Programming a computer for playing chess. \emph{Philosophical Magazine}, 1950.
\bibitem[Tod15]{Tod15} Y.~Todo. Structural evaluation by generalized integral property. \emph{EUROCRYPT 2015}.
\end{thebibliography}

\appendix

\section{Constants}\label{app:const}

\paragraph{S-box $S$} (row: high nibble, column: low nibble, hexadecimal).

% BEGIN SBOX
\begin{center}\small\setlength{\tabcolsep}{3.2pt}
\begin{tabular}{@{}c|cccccccccccccccc@{}}
 & \textbf{\texttt{0}} & \textbf{\texttt{1}} & \textbf{\texttt{2}} & \textbf{\texttt{3}} & \textbf{\texttt{4}} & \textbf{\texttt{5}} & \textbf{\texttt{6}} & \textbf{\texttt{7}} & \textbf{\texttt{8}} & \textbf{\texttt{9}} & \textbf{\texttt{a}} & \textbf{\texttt{b}} & \textbf{\texttt{c}} & \textbf{\texttt{d}} & \textbf{\texttt{e}} & \textbf{\texttt{f}}\\
\hline
\textbf{\texttt{0}} & \texttt{15} & \texttt{86} & \texttt{4d} & \texttt{a2} & \texttt{24} & \texttt{91} & \texttt{fa} & \texttt{cd} & \texttt{75} & \texttt{c1} & \texttt{7e} & \texttt{8d} & \texttt{cb} & \texttt{c8} & \texttt{f6} & \texttt{34}\\
\textbf{\texttt{1}} & \texttt{6d} & \texttt{f5} & \texttt{22} & \texttt{ff} & \texttt{a0} & \texttt{6a} & \texttt{80} & \texttt{12} & \texttt{ab} & \texttt{87} & \texttt{0b} & \texttt{9d} & \texttt{1f} & \texttt{56} & \texttt{5f} & \texttt{df}\\
\textbf{\texttt{2}} & \texttt{2a} & \texttt{84} & \texttt{bf} & \texttt{1b} & \texttt{28} & \texttt{d5} & \texttt{36} & \texttt{30} & \texttt{1e} & \texttt{2f} & \texttt{33} & \texttt{6b} & \texttt{10} & \texttt{68} & \texttt{0c} & \texttt{99}\\
\textbf{\texttt{3}} & \texttt{d3} & \texttt{fb} & \texttt{09} & \texttt{90} & \texttt{0f} & \texttt{ad} & \texttt{aa} & \texttt{42} & \texttt{9c} & \texttt{bd} & \texttt{c3} & \texttt{17} & \texttt{ea} & \texttt{62} & \texttt{38} & \texttt{7b}\\
\textbf{\texttt{4}} & \texttt{45} & \texttt{8c} & \texttt{a9} & \texttt{39} & \texttt{40} & \texttt{83} & \texttt{25} & \texttt{3d} & \texttt{a1} & \texttt{be} & \texttt{b1} & \texttt{41} & \texttt{bb} & \texttt{6e} & \texttt{e6} & \texttt{a8}\\
\textbf{\texttt{5}} & \texttt{d8} & \texttt{95} & \texttt{b8} & \texttt{b7} & \texttt{46} & \texttt{fd} & \texttt{77} & \texttt{5c} & \texttt{e8} & \texttt{14} & \texttt{d7} & \texttt{88} & \texttt{16} & \texttt{4b} & \texttt{13} & \texttt{7f}\\
\textbf{\texttt{6}} & \texttt{ec} & \texttt{81} & \texttt{8f} & \texttt{d9} & \texttt{eb} & \texttt{32} & \texttt{3a} & \texttt{d1} & \texttt{b6} & \texttt{ef} & \texttt{8e} & \texttt{c9} & \texttt{97} & \texttt{9a} & \texttt{50} & \texttt{b9}\\
\textbf{\texttt{7}} & \texttt{ee} & \texttt{72} & \texttt{a4} & \texttt{63} & \texttt{c6} & \texttt{66} & \texttt{f1} & \texttt{00} & \texttt{e1} & \texttt{53} & \texttt{5e} & \texttt{3f} & \texttt{8b} & \texttt{b2} & \texttt{5d} & \texttt{da}\\
\textbf{\texttt{8}} & \texttt{49} & \texttt{57} & \texttt{c4} & \texttt{ce} & \texttt{37} & \texttt{d4} & \texttt{cf} & \texttt{af} & \texttt{b4} & \texttt{19} & \texttt{b0} & \texttt{e3} & \texttt{31} & \texttt{01} & \texttt{76} & \texttt{26}\\
\textbf{\texttt{9}} & \texttt{5a} & \texttt{ba} & \texttt{04} & \texttt{03} & \texttt{70} & \texttt{f4} & \texttt{f0} & \texttt{47} & \texttt{f9} & \texttt{f8} & \texttt{8a} & \texttt{e0} & \texttt{35} & \texttt{4e} & \texttt{73} & \texttt{27}\\
\textbf{\texttt{a}} & \texttt{21} & \texttt{43} & \texttt{c7} & \texttt{0a} & \texttt{59} & \texttt{e5} & \texttt{2d} & \texttt{78} & \texttt{e9} & \texttt{a3} & \texttt{d2} & \texttt{6c} & \texttt{89} & \texttt{e2} & \texttt{4a} & \texttt{94}\\
\textbf{\texttt{b}} & \texttt{05} & \texttt{5b} & \texttt{51} & \texttt{cc} & \texttt{55} & \texttt{fe} & \texttt{7d} & \texttt{c0} & \texttt{b5} & \texttt{ca} & \texttt{79} & \texttt{93} & \texttt{48} & \texttt{db} & \texttt{e4} & \texttt{3c}\\
\textbf{\texttt{c}} & \texttt{bc} & \texttt{ed} & \texttt{54} & \texttt{9f} & \texttt{64} & \texttt{60} & \texttt{65} & \texttt{82} & \texttt{a5} & \texttt{7a} & \texttt{02} & \texttt{2b} & \texttt{6f} & \texttt{98} & \texttt{08} & \texttt{71}\\
\textbf{\texttt{d}} & \texttt{69} & \texttt{ac} & \texttt{1c} & \texttt{07} & \texttt{20} & \texttt{f7} & \texttt{4f} & \texttt{61} & \texttt{23} & \texttt{67} & \texttt{b3} & \texttt{f3} & \texttt{9e} & \texttt{a6} & \texttt{1d} & \texttt{29}\\
\textbf{\texttt{e}} & \texttt{58} & \texttt{96} & \texttt{74} & \texttt{f2} & \texttt{a7} & \texttt{11} & \texttt{1a} & \texttt{9b} & \texttt{3b} & \texttt{52} & \texttt{e7} & \texttt{dc} & \texttt{06} & \texttt{85} & \texttt{c2} & \texttt{4c}\\
\textbf{\texttt{f}} & \texttt{2c} & \texttt{44} & \texttt{3e} & \texttt{2e} & \texttt{ae} & \texttt{7c} & \texttt{92} & \texttt{de} & \texttt{dd} & \texttt{0d} & \texttt{d0} & \texttt{18} & \texttt{0e} & \texttt{fc} & \texttt{c5} & \texttt{d6}\\
\end{tabular}
\end{center}
% END SBOX

\paragraph{Matrices $M$ ($\theta$) and $M^{-1}$ ($\theta^{-1}$)} (hexadecimal, row by row).

% BEGIN MDS
\begin{center}\small\setlength{\tabcolsep}{4pt}
\begin{tabular}{@{}cccccccc@{}}
\multicolumn{8}{c}{$M$}\\[2pt]
\texttt{6a} & \texttt{d7} & \texttt{84} & \texttt{ac} & \texttt{af} & \texttt{d6} & \texttt{e3} & \texttt{61}\\
\texttt{de} & \texttt{e3} & \texttt{72} & \texttt{e6} & \texttt{d3} & \texttt{eb} & \texttt{d7} & \texttt{17}\\
\texttt{6d} & \texttt{08} & \texttt{d8} & \texttt{1b} & \texttt{16} & \texttt{b1} & \texttt{4e} & \texttt{51}\\
\texttt{32} & \texttt{4e} & \texttt{8a} & \texttt{fc} & \texttt{5e} & \texttt{0d} & \texttt{08} & \texttt{ec}\\
\texttt{8a} & \texttt{1e} & \texttt{32} & \texttt{53} & \texttt{36} & \texttt{cf} & \texttt{b3} & \texttt{22}\\
\texttt{d8} & \texttt{b3} & \texttt{6d} & \texttt{04} & \texttt{43} & \texttt{ad} & \texttt{1e} & \texttt{f0}\\
\texttt{72} & \texttt{5d} & \texttt{de} & \texttt{a9} & \texttt{49} & \texttt{c6} & \texttt{50} & \texttt{18}\\
\texttt{84} & \texttt{50} & \texttt{6a} & \texttt{27} & \texttt{a6} & \texttt{0e} & \texttt{5d} & \texttt{3e}\\
\end{tabular}
\qquad
\begin{tabular}{@{}cccccccc@{}}
\multicolumn{8}{c}{$M^{-1}$}\\[2pt]
\texttt{f1} & \texttt{a7} & \texttt{cf} & \texttt{26} & \texttt{40} & \texttt{46} & \texttt{f8} & \texttt{48}\\
\texttt{36} & \texttt{69} & \texttt{5f} & \texttt{5b} & \texttt{d4} & \texttt{15} & \texttt{11} & \texttt{c7}\\
\texttt{4d} & \texttt{b9} & \texttt{41} & \texttt{23} & \texttt{50} & \texttt{a3} & \texttt{55} & \texttt{92}\\
\texttt{80} & \texttt{53} & \texttt{52} & \texttt{4d} & \texttt{f4} & \texttt{8b} & \texttt{31} & \texttt{19}\\
\texttt{ed} & \texttt{73} & \texttt{15} & \texttt{f2} & \texttt{78} & \texttt{4f} & \texttt{aa} & \texttt{26}\\
\texttt{99} & \texttt{ae} & \texttt{e0} & \texttt{09} & \texttt{6f} & \texttt{d6} & \texttt{1e} & \texttt{68}\\
\texttt{e6} & \texttt{e2} & \texttt{05} & \texttt{06} & \texttt{d1} & \texttt{2e} & \texttt{74} & \texttt{a8}\\
\texttt{2b} & \texttt{85} & \texttt{a1} & \texttt{97} & \texttt{df} & \texttt{11} & \texttt{f5} & \texttt{3e}\\
\end{tabular}
\end{center}
% END MDS

\paragraph{Padding constant} (bits $410..639$ of $L$, as an integer):
\[
  \mathit{pad} = \texttt{0x391d69042f080a4aa81f1dc9676116fafcd238b21e43a01e5fb06f0fe6}.
\]

\section{Test vectors (v1.1, hexadecimal)}\label{app:vectors}

% BEGIN VECTORS
\begin{center}\small
\begin{tabular}{@{}ll@{}}
\toprule
\multicolumn{2}{@{}l}{\emph{Vector 1}}\\
Key & $\hx{00}^{52}$ (52 zero bytes)\\
Tweak & $\hx{00}^{16}$\\
Plaintext & $\hx{00}^{64}$\\
Ciphertext & \texttt{67b7e691273da3b40bedf3cf4abad7a3c960473b1ef4356d961413423c718053}\\ & \texttt{2fab02454f11cb600955edd094de6490024eb8f06f40d4f0d30443e004069d72}\\
\midrule
\multicolumn{2}{@{}l}{\emph{Vector 2}}\\
Key & bytes $0..50$ equal to $0,1,\dots,50$ ($\hx{00}\,\hx{01}\cdots\hx{32}$), byte $51 = \hx{02}$\\
Tweak & \texttt{000102030405060708090a0b0c0d0e0f}\\
Plaintext & bytes $0..63$ equal to $0,1,\dots,63$ ($\hx{00}\,\hx{01}\cdots\hx{3f}$)\\
Ciphertext & \texttt{9f52662ad01ba4bafe6fa894e81fd6b23fa22078bb5054587d09214cd8306ba1}\\ & \texttt{c3bf6167b07494d2d5f379c26b991a79403de492f4e310639917d33af96ada4c}\\
\bottomrule
\end{tabular}
\end{center}
% END VECTORS

Twelve block-cipher vectors (with $\RK_0$ and $\RK_{16}$) and six BRI-410 AE vectors are in
\texttt{publikim/reference/test\_vectors.json}.

\section{The CP-SAT model}\label{app:cpsat}

For rounds $t = 0,\dots,R$ and cells $c$, the binary variable $x_{t,c} \in \{0,1\}$ states that cell $c$ is active
at the input of the S-layer of round $t+1$. For every round $t$ and row $q$, a binary variable $d_{t,q}$ marks the
row as active, and every active input or output cell of the row implies $d_{t,q}$. A binary variable $s_t$ selects
$d_r = +1$ or $d_r = -1$ in round $t$, so the solver chooses the worst case. The inputs of row $q$ of $\theta$ are
the cells $(r,k)$ with $r + k\,d_r \equiv q \pmod 8$, and the constraint is
\[
  \sum_{(r,k):\; r + k d_r \equiv q} x_{t,(r,k)} \;+\; \sum_{k=0}^{7} x_{t+1,(q,k)} \;\ge\; 9\, d_{t,q} .
\]
The objective is
\[
  \min \sum_{t=0}^{R-1} \sum_{c} x_{t,c} \qquad \text{subject to} \qquad \sum_c x_{0,c} \ge 1 .
\]
The model of the schedule is the same, with 10 rows and the map $j \mapsto 9j \bmod 80$, and it starts from a
non-zero difference in bytes $0..15$ (the tweak). Code: \texttt{matje/f4b\_milp.py}.

\end{document}
